% Patch record for C-62 -- humanizer pass over Sebas's round-2 \revblue{} text (academic register).
% Applied 2026-10-29 (Sebas). No number, citation or claim changed.
% results.tex l.945 also carries the C-61 LiDAR fix in the same line.
% validate_tex.sh OK; check_roundtrip.sh OK; compile_pdf.sh OK.
%
% --- a/sections/conclusions.tex
% +++ b/sections/conclusions.tex
% @@ -16 +16 @@ Section~\ref{ssec:Limitations} details the present work's limitations, spanning
% -\revblue{The findings show that lightweight, bio-inspired neural networks ---arbitrating between locomotion modes and reactive behaviors driven by visual and inertial cues--- can manage multimodal gait transitions on a real robotic platform, with the neural circuits executed on an external computer.} The perception-decision pipeline additionally retained task-level performance under combined IMU-drift and LiDAR-noise perturbation, and, conditionally on the scenario, camera illumination distortion (Section~\ref{ssec:NoiseRobustnessEval}), evaluated jointly rather than as isolated single-sensor degradations. \revblue{Although the physical prototype displays hardware-level boundaries (Section~\ref{ssec:Limitations}),} these results substantiate that complex, multi-stimulus mobility decisions can be \revblue{taken for such resource-constrained platforms; onboard feasibility was assessed only for the MLP terrain classifier on the ESP32-S3 (Section~\ref{sssec:mlp_eval}), not for the full neural pipeline}.
% +\revblue{The findings show that lightweight, bio-inspired neural networks that arbitrate between locomotion modes and reactive behaviors from visual and inertial cues can manage multimodal gait transitions on a real robotic platform, with the neural circuits executed on an external computer.} The perception-decision pipeline additionally retained task-level performance under combined IMU-drift and LiDAR-noise perturbation, and, conditionally on the scenario, camera illumination distortion (Section~\ref{ssec:NoiseRobustnessEval}), evaluated jointly rather than as isolated single-sensor degradations. \revblue{Although the physical prototype displays hardware-level boundaries (Section~\ref{ssec:Limitations}),} these results substantiate that complex, multi-stimulus mobility decisions can be \revblue{taken for such resource-constrained platforms; onboard feasibility was assessed only for the MLP terrain classifier on the ESP32-S3 (Section~\ref{sssec:mlp_eval}), not for the full neural pipeline}.
% --- a/sections/introduction.tex
% +++ b/sections/introduction.tex
% @@ -10 +10 @@ Wheeled-legged robots address this limitation by integrating actuated limbs (pro
% -A critical observation emerges from this body of work: in all published wheeled-legged systems, locomotion mode switching is resolved either heuristically through hand-coded threshold logic \citep{Kashiri2019CENTAURO}, by model predictive control \citep{bjelonic2019,Bjelonic2022OfflineMotion}, or by end-to-end RL policies that encode switching implicitly as a black-box behavior \citep{Lee2024Science,Humphreys2025BioInspired}. None of these approaches employs a neurobiologically interpretable arbitration architecture. \revblue{Table~\ref{table:action_selection_comparison} contrasts these approaches with central pattern generators (CPGs) ---including feedback-modulated variants that trigger gait transitions from vestibular or load feedback--- and with the basal-ganglia winner-take-all arbitration adopted here.}
% +A critical observation emerges from this body of work: in all published wheeled-legged systems, locomotion mode switching is resolved either heuristically through hand-coded threshold logic \citep{Kashiri2019CENTAURO}, by model predictive control \citep{bjelonic2019,Bjelonic2022OfflineMotion}, or by end-to-end RL policies that encode switching implicitly as a black-box behavior \citep{Lee2024Science,Humphreys2025BioInspired}. None of these approaches employs a neurobiologically interpretable arbitration architecture. \revblue{Table~\ref{table:action_selection_comparison} contrasts these approaches with central pattern generators (CPGs), including feedback-modulated variants that trigger gait transitions from vestibular or load feedback, and with the basal-ganglia winner-take-all arbitration adopted here.}
% --- a/sections/results.tex
% +++ b/sections/results.tex
% @@ -480 +481 @@ locomotion phase.
% -\revblue{Table~\ref{tab:fsm_comparison} compares the FSM controller of Section~\ref{ssec:FSMBaseline} with the basal-ganglia network on all seven simulated suites, reporting success rate, task completion time, Roll RMS, and $\lambda$ --- a motor-command clarity ratio (commanded-velocity magnitude relative to neural activity norm; higher indicates more decisive output).}
% +\revblue{Table~\ref{tab:fsm_comparison} compares the FSM controller of Section~\ref{ssec:FSMBaseline} with the basal-ganglia network on all seven simulated suites, reporting success rate, task completion time, Roll RMS, and $\lambda$, a motor-command clarity ratio (commanded-velocity magnitude relative to neural activity norm; higher indicates more decisive output).}
% @@ -822 +823 @@ To further consolidate the physical validation under complex, multi-section topo
% -\revblue{The arbitration network autonomously transitions through this sequence.} The critical event occurs in Frames 5--6: rolling-mode wheel slip on the stochastic rough segment triggers a sustained stagnation signal, resolved by reactively redeploying the quadrupedal gait, which restores the support polygon and completes traversal. This closed-loop reactive switching was observed across all 15 trials, demonstrating reliable, autonomous failure detection and gait recovery under stochastic, non-structured terrain conditions.
% +\revblue{The arbitration network runs through this sequence without operator input.} The critical event occurs in Frames 5--6: rolling-mode wheel slip on the stochastic rough segment triggers a sustained stagnation signal, resolved by reactively redeploying the quadrupedal gait, which restores the support polygon and completes traversal. This closed-loop reactive switching was observed across all 15 trials, demonstrating reliable, autonomous failure detection and gait recovery under stochastic, non-structured terrain conditions.
% @@ -944 +945 @@ To further consolidate the physical validation under complex, multi-section topo
% -\textbf{Complex (see Figure~\ref{fig:complex_real_compound}):} \revblue{Across successive back-to-back transformations, the deterministic step response shows the ROS~2 Control architecture managing electromechanical inertia and trajectory blending, protecting the joints from cumulative torque stress.} The three-phase firing-variance profile maps the multi-stage decision process: a transient hyper-excitation peak during static obstacle avoidance, driven by the physical barrier's high spatial density flooding the camera sensor; a secondary activation wave capturing the emergency flight reflex to the aversive cue; and a low-energy steady-state as the robot tracks the appetitive target to mission completion. The chassis attitude telemetry mirrors this sequence, with asymmetric flanking maneuvers exciting the transversal axis and rapid backward accelerations peaking the longitudinal tilt, before both metrics drop and stabilize during the final convergence phase, indicating the wheel-driven configuration smooths parasitic oscillations during target approach.
% +\textbf{Complex (see Figure~\ref{fig:complex_real_compound}):} \revblue{The step response stays deterministic across back-to-back transformations: the ROS~2 Control layer absorbs the electromechanical inertia and blends the trajectories, so torque stress does not accumulate in the joints.} The three-phase firing-variance profile maps the multi-stage decision process: a transient hyper-excitation peak during static obstacle avoidance, \revblue{driven by the dense LiDAR returns from the physical barrier}; a secondary activation wave capturing the emergency flight reflex to the aversive cue; and a low-energy steady-state as the robot tracks the appetitive target to mission completion. The chassis attitude telemetry mirrors this sequence, with asymmetric flanking maneuvers exciting the transversal axis and rapid backward accelerations peaking the longitudinal tilt, before both metrics drop and stabilize during the final convergence phase, indicating the wheel-driven configuration smooths parasitic oscillations during target approach.
% @@ -1034 +1035 @@ To evaluate repeatability and determinism under uncalibrated physical environmen
% -As an open, low-cost exploratory research baseline, the physical robot is constructed using low-torque servomotors and a lightweight PLA structure, which limits its operational envelope compared to industrial-grade platforms. Locomotion gaits are currently executed via pre-configured open-loop control patterns; while this reduces control complexity, it weakens active stabilization capabilities on high-slip surfaces or excessively steep inclinations. \revblue{This constraint stems from the hardware, not from the neural arbitration architecture.} As a representative case, the gait selected during steep-slope negotiation (Figure~\ref{fig:RES_IMU2_real}) reflects the boundaries of the physical actuator envelope: the neural decision module consistently selects the locomotion mode associated with improved stability under the detected terrain conditions, yet the open-loop execution of that gait cannot compensate for the torque deficits inherent to the hardware. Experimental validation was constrained to structured laboratory environments with static stimuli and did not include unconstrained outdoor topographies, dynamic obstacles, or severe external disturbances. Therefore, the empirical results presented herein do not constitute validation of a production-ready system for rugged industrial deployment. Instead, they substantiate the mathematical, biological, and computational feasibility of a highly efficient basal-ganglia inspired arbitration network capable of managing multimodal mobility \revblue{on a low-cost robotic platform}.
% +As an open, low-cost exploratory research baseline, the physical robot is constructed using low-torque servomotors and a lightweight PLA structure, which limits its operational envelope compared to industrial-grade platforms. Locomotion gaits are currently executed via pre-configured open-loop control patterns; while this reduces control complexity, it weakens active stabilization capabilities on high-slip surfaces or excessively steep inclinations. \revblue{This limit lies in the hardware and is independent of the neural arbitration architecture.} As a representative case, the gait selected during steep-slope negotiation (Figure~\ref{fig:RES_IMU2_real}) reflects the boundaries of the physical actuator envelope: the neural decision module consistently selects the locomotion mode associated with improved stability under the detected terrain conditions, yet the open-loop execution of that gait cannot compensate for the torque deficits inherent to the hardware. Experimental validation was constrained to structured laboratory environments with static stimuli and did not include unconstrained outdoor topographies, dynamic obstacles, or severe external disturbances. Therefore, the empirical results presented herein do not constitute validation of a production-ready system for rugged industrial deployment. Instead, they substantiate the mathematical, biological, and computational feasibility of a highly efficient basal-ganglia inspired arbitration network capable of managing multimodal mobility \revblue{on a low-cost robotic platform}.
